\documentclass[12pt]{article}
\usepackage[utf8]{inputenc}
\usepackage[T1]{fontenc}
\usepackage{hyperref}
\title{Toraja-Sa'dan Language Resources for K-12 Curriculum Development}
\author{Compiled by Claude Sonnet 5.5 (claude-sonnet-5-5) for Andr\'as Kornai}
\date{}
\begin{document}
\maketitle

\section{Introduction}
Toraja-Sa'dan (ISO 639-3 sda; Glottocode tora1261) is an Austronesian language of South Sulawesi, Indonesia. Speakers call it \emph{basa Toraya}; older Dutch sources call it Tae'. Population estimates vary (about 590,000 in the 2010 census according to Wikipedia; about 750,000 in the survey's reference description). Searches found little in the way of published school materials, but they did find a historical dictionary, a Bible translation, and several academic studies of the local-content (\emph{muatan lokal}) subject ``Bahasa Toraja''. This report lists only what appeared in search results.

\section{Primary Educational Resources}
Search results did not turn up an official textbook or primer for Bahasa Toraja. They did turn up research on teaching it as a local-content subject.
\begin{itemize}
\item Eprints record of a development study of interactive digital multimedia for the Bahasa Toraja local-content subject at SMA Kristen Makale (grade X), 2024: \url{https://eprints.unm.ac.id/36141/}
\item ``Need analysis of learning local content of bahasa Toraja in elementary school students'' (Kembara journal, 2022): \url{https://ejournal.umm.ac.id/index.php/kembara/article/view/21940}
\item Study of a Toraja language-variation lesson for primary pupils at Ke'te' Kesu': \url{https://ejournal.undiksha.ac.id/index.php/JPI/article/download/28622/21767/120397}
\item Tongkonan, a Ministry of Education children's reader on Toraja culture (Weni Rahayu, 2017). It is not in Toraja but is relevant as culture content: \url{https://lobo.apps01.yorku.ca/wp-content/uploads/2021/09/Rahayu-2017.pdf}
\end{itemize}

\section{Linguistic and Reference Materials}
\begin{itemize}
\item Wikipedia overview of the language: \url{https://www.Wikipedia.com/wiki/Toraja_language}
\item Omniglot page on Toraja-Sa'dan: \url{https://omniglot.com/writing/torajasadan.htm}
\item The principal dictionary is the Toraja-Dutch dictionary of H. van der Veen (1940), with an Indonesian translation in 1972, according to the search summary. No online copy was found in the search results.
\item Study on the design of an Android Toraja dictionary app, which notes the printed dictionary circulated only in small numbers: \url{https://ejournal.upbatam.ac.id/index.php/comasiejournal/article/view/6950}
\item Glottolog and Wikidata entries for identification: \url{https://www.wikidata.org/wiki/Q36673}
\end{itemize}

\section{Religious and Community Resources}
\begin{itemize}
\item Scripture Earth page for Toraja (sda), listing a 2009 Indonesian Bible Society (LAI) translation, a children's version audio and a Gospel of Mark film: \url{https://www.scriptureearth.org/00eng.php?iso=sda}
\item LAI appeal for 5,000 illustrated Good News booklets for children in Toraja: \url{https://www.alkitab.or.id/en/berita/2166-lets-support-the-publication-of-5000-illustrated-good-news-in-toraja-language}
\item Joshua Project profile of the Toraja-Sa'dan people: \url{https://joshuaproject.net/index.php/people_groups/15565}
\end{itemize}

\section{Digital and Audio Resources}
\begin{itemize}
\item The Scripture Earth page above lists audio and video Scripture items.
\item The Android dictionary project above is a software lead, but no released app was confirmed.
\end{itemize}

\section{Curriculum Development Context}
Although the reference description says the language is not known to be taught in schools, the studies above describe Bahasa Toraja as a local-content subject in primary and secondary schools in Tana Toraja. How widely this is implemented, and what official materials exist, was not established; the Tana Toraja education office would need to be consulted. Concrete gaps are a modern open dictionary, graded primers, and a standard orthography guide. Bible translations and the van der Veen dictionary could seed a corpus.
\end{document}